\begin{statement}{C.2.UN.1}{Proved}{Every torus-degeneration fibre has even Euler number}
Let $f:Y\to B$ be a proper holomorphic map from a smooth complex
threefold to a smooth complex curve, with connected fibres and smooth
general fibre a complex two-torus. For every $b\in B$,
\[
                    e((Y_b)_{\rm red})\in2\Z.       \tag{C.UN.1}
\]
This assertion allows arbitrary multiplicities and singularities and
imposes no K\"ahler or Betti-number hypothesis on $Y$.
\end{statement}
\begin{proof}
Take a sufficiently small closed coordinate disc $\Delta$ about $b$,
containing no other critical value. Put $U=f^{-1}(\Delta)$ and
$M=\partial U$. Then $U$ is a compact oriented smooth real six-manifold
with boundary, and $M$ is the mapping torus of the orientation-preserving
monodromy of a real four-torus. The analytic central support $S$ is a
deformation retract of $U$. One precise construction uses a controlled
gradient for the real analytic function $|f|^2$ on a Whitney
stratification: it decreases $|f|$, is tangent to the strata over zero,
and has trajectories with a limit on the zero set by the local
analytic gradient inequality. Compactness of the zero set gives a
common small disc, and the resulting continuous limit homotopy fixes
the zero set. Equivalently the triangulated proper local tube is a
regular neighbourhood of the zero fibre. Nilpotents do not change this
retraction. Thus $e(S)=\chi(U)$.

We give the parity calculation entirely over $\Q$. Write
$s_i=\operatorname{rank}(H_i(U)\to H_i(U,M))$.
Exactness of the pair sequence, with Poincar\'e--Lefschetz duality,
gives
\[
 b_0(M)+b_1(M)+b_2(M)
       =\sum_{i=0}^{5}b_i(U)-s_0-s_3-2s_1-2s_2.
                                                        \tag{C.UN.2}
\]
Indeed, the summand $b_i(M)$ is the sum of the ranks of
$H_i(M)\to H_i(U)$ and $H_{i+1}(U,M)\to H_i(M)$;
substitution of $b_{i+1}(U,M)=b_{5-i}(U)$ gives the formula.
Here $s_0=0$ and $b_6(U)=0$. The map defining $s_3$ is represented,
after duality, by the middle-dimensional intersection pairing of $U$.
That pairing is skew-symmetric because the middle dimension is three,
so $s_3$ is even. Consequently
\[
 \chi(U)\equiv b_0(M)+b_1(M)+b_2(M)\pmod2.            \tag{C.UN.3}
\]

Let $T$ act on $H^1(F,\Q)$ for the four-torus fibre $F$, and set
$a_i=\dim\ker(\Lambda^iT-I)$. Cutting the mapping torus along one
fibre gives the Wang exact sequence; hence
\[
 b_0(M)=1,\quad b_1(M)=1+a_1,\quad b_2(M)=a_1+a_2.
                                                        \tag{C.UN.4}
\]
Orientation preservation says $\det T=1$. The transformation
$G=\Lambda^2T$ preserves the nondegenerate symmetric wedge pairing
on $\Lambda^2\Q^4$ and has determinant $(\det T)^3=1$.
We verify that $\operatorname{rank}(G-I)$ is even without a
semisimplicity assumption. Complexify and split off the generalized
$(-1)$-eigenspace. Its dimension is even: all eigenvalues outside
$\{1,-1\}$ occur in reciprocal pairs with equal algebraic
multiplicity, and $\det G=1$. This generalized eigenspace is
nondegenerate and orthogonal to its invariant complement. On it
$G-I$ is invertible, of even rank. On the complement $G+I$ is
invertible, and
\[
                         (G-I)(G+I)^{-1}
\]
is skew-adjoint for the symmetric wedge pairing. Multiplication by
that nondegenerate pairing represents it as a skew-symmetric matrix,
whose rank is even. Its rank equals that of $G-I$ on the complement.
This proves the assertion. Since $\dim\Lambda^2\Q^4=6$, $a_2$ is
even. Equations (C.UN.3)--(C.UN.4) prove (C.UN.1).
\end{proof}

\begin{statement}{C.2.UN.2}{Proved}{All even Euler numbers and an unbounded zero-Euler family}
The set of Euler numbers of reduced fibres among compact smooth
threefolds fibred in two-tori is exactly $2\Z$. For each $k\ge1$
there is a smooth projective such threefold with one singular reduced
normal-crossing fibre of Euler number zero consisting of a blown-up
abelian surface, $k$ ruled surfaces over a genus-two curve, and $k$
projective planes. Their incidence graph is a chain of the ruled
components attached to the abelian component, together with $k$
projective-plane leaves attached to that component.
\end{statement}
\begin{proof}
Use the explicit embedded genus-two curve $C_2\subset A$ of C2b.
In $A\times\PP^1$ retain the smooth horizontal surface
$R=C_2\times\PP^1$. First blow up $C_2\times\{0\}$, as in C2b.
The strict transform of $R$ is still $R$, because its centre is a
Cartier divisor in $R$. It meets the exceptional ruled surface in
the section corresponding to the base-normal direction; this
section is disjoint from the intersection with the strict abelian
component. The map to the base is a submersion near this section.

Blow up that section and repeat, always using the intersection of
the strict transform of $R$ with the newest exceptional component.
The next centre is a smooth genus-two curve in the smooth locus of
the reduced central fibre, disjoint from its double curves. Its
normal sequence splits: the normal line to the curve inside $R$
maps isomorphically under $df$ to the tangent line of the base,
and supplies a splitting of the normal sequence of the central
surface. The two line summands give two disjoint sections on the
new exceptional projective bundle. One meets the previous
component; the other is the next centre. This proves by induction
that all $k$ curve blowups are defined and that the reduced central
support has the stated chain. In local coordinates the comparison
is always $t=u\tau$ as in C2b, so every new component has
multiplicity one and no triple point is introduced.

Finally blow up $m\ge0$ distinct points of the strict abelian
component away from the genus-two double curve. They introduce
$m$ projective-plane leaves, each meeting the blown-up abelian
component in one projective line. All maps are proper holomorphic
blowdowns followed by the original projection. Every nonzero
fibre is biholomorphic to $A$.

Each curve blowup increases the Euler number of the central
support by $e(C_2)=-2$: the exceptional projective bundle has Euler
number $2e(C_2)$ and its intersection with the old component has
Euler number $e(C_2)$. Each point blowup increases it by two,
because the old surface gains one, the exceptional plane adds
three, and their intersection subtracts two. Thus
\[
                         e(S)=2(m-k).                 \tag{C.UN.5}
\]
Given any integer $n$, take $k\ge\max(1,-n)$ and $m=k+n$.
This realizes $2n$; C.2.UN.1 supplies the opposite inclusion.
Taking $m=k$ gives the asserted zero-Euler family.

The integral cohomology of the total space is torsion-free. The
excision and projective-bundle comparison of C2b gives
\[
 (b_0,\ldots,b_6)
      =(1,4,7+k+m,8+4k,7+k+m,4,1).                  \tag{C.UN.6}
\]
Its Euler characteristic is independently $2(m-k)$. The structure
sheaf direct images of each blowup are $\OO$ in degree zero and
zero in positive degrees, so
$(h^0(\OO),h^1(\OO),h^2(\OO),h^3(\OO))=(1,2,1,0)$.
The exact arithmetic and exterior-square checks are in
\path{code/root_unrestricted_euler.py} and its output.
This family supplies further members of the requested unrestricted
zero-Euler classification; it does not claim that every zero-Euler
fibre has this construction.
\end{proof}
